\chapter{Diskussion}
\label{ch:Diskussion}

\section{Zusammenfassung}
\label{sec:Zusammenfassung}
Das Ziel des Versuches war es, die Viskosität von PEG zu bestimmen, sowie zu schauen, ob die jeweilige Nährung von laminarer Strömung Gültigkeit hat.
Die Bedetung der Messergebnisse soll in \csec{Analyse} diskutiert werden.

\subsection{Stokes}
In diesem Versuchsteil sollte, wie erwähnt, die Viskosität des PEG bestimmt werden, die theoretische Sinkgeschwindigkeit bestimmt werden und geschaut werden, für welche Kugelradien die Nährung für eine laminare Strömung gültig ist. 

\paragraph{Viskosität}
Zur Bestimmung der Viskosität wurde die mittlere Sinkgeschwindigkeit der einzelnen Kurgeln bestimmt. Zudem musste die Dichte des PEG anhand seiner Temperatur bestimmt werden. Somit waren alle Werte zur Bestimmung zusammen.

Über die Messergebnisse konnte die Viskosität von 
\res{\eta_\text{s}}{0.263}{0.017}{\frac{kg}{m\, s}}
bestimmt werden. 

\paragraph{theoretische Sinkgeschwindigkeit}
Die Messergebnisse der Sinkgeschwindigkeit wurden mit der Ladenburg-Korrektur nachverbessert. Später wurde dann die theoretisch Sinkgeschwindigkeit in laminarer Strömung bestimmt.
Die Ergebnisse sind tabellarisch im \hyperref[ch:Protokoll]{Messprotokoll} und \ctab{reynolds_tab} festgehalten.

\paragraph{Reynoldszahl}
Für die Reynodls-Zahl wurde der Plot in \cabb{plot/ver_ges_gegen_re} erstellt. Dieser hatte nicht die erhoffte Form. Der kritische Wert musste daher unter Annahme eines großen Fehlers abgelesen werden:

\res{Re_\text{krit}}{0.6}{0.3}{}

Dabei wurde festgestellt, das die Ungenauigkeit der Temperatur keine besonders wichtige Rolle spielt und der Wert von $\Delta T = \pm 0.5  \mathrm{^\circ C}$ angenommen werden kann.

\subsection{Hagen-Poiseuille}
Auch in diesem Versuchsteil sollte die Viskosität des PEG bestimmt werden, sowie die Gültigkeit der laminaren Nährung geprüft werden.

\paragraph{Viskosität}
Die Viskosität nach Hagen-Poiseuille wurde über den Volumenfluss einer Kapillare bekannter Maße bestimmt. Dafür wurden die Messwerte geplottet und die Steigung bestimmt. Zuletzt wurde der Druck über die Höhendifferenz vor und nach der Zeitstoppung bestimmt. Aus allen Werten konnte auf die Viskosität 

\res{\eta_\text{hp}}{0.222}{0.006}{\frac{kg}{m\, s}}

geschlossen werden.

\paragraph{Reynoldszahl}
Über die durchschnittliche Ausflussgeschwindigkeit konnte auch für dieses Verfahren die Reynodls-Zahl bestimmt werden. Bei dieser Methode kam ein Wert von 

\res{Re_\text{hp}}{0.1617}{0.0016}{ }

heraus. Dass die Reynoldszahelen nicht zusammenzupassen scheinen ist ganz richtig, denn die Reynodlszahl muss für jedes System seperat betrachtet werden. 

\section{Analyse der Messwerte}
\label{sec:Analyse}
Nun sollen die Werte noch analysiert werden und überprüft werden, ob die Werte deckungsgleich sind. 

\paragraph{Viskosität} Zunächst soll geschaut werden, ob die Werte der Viskosität sich statistisch decken können. Dafür wird die Standardabweichung genutzt. Nach der Berechnung kommt man auf eine Abweichung von

\eq{sig_visk}{
    \sigma_\text{visk} = 2.27\sigma
}

Damit lässt sich sagen, dass die Ergebnisse zwar noch in einer $3\sigma$-Umgegung liegen und somit im statistisch signifikanten Zusammenhang stehen. Jedoch wären bessere Ergebnisse natürlcih wünschenswert gewesen. 
Da es leider keinen Literaturwert gibt, können wir hier nicht sagen, welche Methode wirklich die bessere Messmethode darstellt. 

Jedoch ist in der ersten Messmethode die menschliche Ungenauigkeit viel mehr von bedeutung, als in der nach Hagen-Poiseuille. Denn in der ersten Methode sind die Zeitintervalle kleiner gewesen, wodurch die Reaktionszeit schwerer wiegt. Zudem war es fast unmöglich, die Kugeln präzise zentrisch ins Sinkrohr einzulassen, dadurch waren die Zeiten noch unpräziser. Zum ausgleich waren fünf Messungen vermutlich nicht genug.

Aber auch die Messmethode nach Hagen-Poiseuille behart auf menschliche Reaktionen. Hier ist die Reaktionszeit jedoch weniger massgeblich. Jedoch ist das Ablesen des Füllstandes eine Herausforderung für das menschliche Auge, da die Oberfläche eine starke Krümmung aufweist und hier das Ablesen vermtulich schwierikger ist, als in der Methode nach Stokes. Auch muss ein konstanter Massefluss sichergestellt sein. Läuft hier etwas schief, so verschieben sich die Werte. So ist es nicht auszuschließen, dass die POM-Kugeln am Boden zu einer leichen Verstopfung füren könnten. Diese Vermutung legt die \cabb{Versuchsaufbau1} nahe, da hier unten bereits sehr viele Kugeln liegen und es visuell so scheint, als leigen diese vor der Öffnung.

Dennoch wird die These hier aufgestellt, dass die Methode nach Hagen-Poiseuille genauer ist, wie im Paragraphen \afz{Reynoldszahlen} erläutert werden soll.

\paragraph{Geschwindigkeiten}

Doch zunächst sollen noch die Geschwindigkeiten diskutiert werden. Wie hilfreich war die Ladenburg-korrektur und wie weit weichen selbst die korrigierten Sinkgeschwindigkeiten von den theoretischen Werten ab?
Die Antworten auf diese Fragen sind tabellarisch bantwortet (\ctab{disk_ges}).

\begin{table}[h]
    \caption{Vergleich und Bestimmugn der Signifikanz der verschiedenen Geschwindigkeiten für die Kugel}
    \label{tab:disk_ges}
    \centering
    \begin{tabular}{L L L || L | L }
        \toprule
        v_\text{gem} [\mathrm{m/s}] & v_\text{korr} [\mathrm{m/s}] & v_\text{theo} [\mathrm{m/s}] & \sigma_\text{gem, korr} & \sigma_\text{korr, theo} \\ 
        \midrule
        0.63 \pm 0.03   & 0.641 \pm 0.026   & 0.48 \pm 0.05     & 0.28 \sigma & 2.86 \sigma \\
        2.00 \pm 0.08   & 2.06 \pm 0.09     & 1.91 \pm 0.18     & 0.50 \sigma & 0.75 \sigma \\
        4.55 \pm 0.07   & 4.74 \pm 0.11     & 4.2 \pm 0.4       & 1.46 \sigma & 1.30 \sigma \\
        7.48 \pm 0.07   & 7.90 \pm 0.16     & 7.4 \pm 0.7       & 2.40 \sigma & 0.70 \sigma \\
        11.63 \pm 0.23  & 12.4 \pm 0.3      & 11.5 \pm 1.0      & 2.04 \sigma & 0.86 \sigma \\
        15.02 \pm 0.20  & 16.3 \pm 0.4      & 16.3 \pm 1.5      & 2.86 \sigma & 0.00 \sigma \\
        20.53 \pm 0.18  & 22.5 \pm 0.7      & 21.90 \pm 1.10    & 2.73 \sigma & 0.46 \sigma \\
        25.3 \pm 0.4    & 28.15 \pm 1.11    & 28.2 \pm 2.6      & 2.42 \sigma & 0.02 \sigma \\
        30.9 \pm 0.3    & 34.8 \pm 1.6      & 35 \pm 3          & 2.40 \sigma & 0.06 \sigma \\
        \bottomrule
    \end{tabular}
\end{table}

Die Werte zeigen zum einen sehr schön die Beobachtung, die in \cabb{plot/ges_gegen_r2} gemacht wurde, denn die Ladenburg-Korrektur macht sich besonders bei größeren Radien bemerkbar. 
Zudem wird deutlich, dass selbst die Ladenburg-Korrektur zwar nicht vollständig ausreicht, um die turbulente Strömung wegzukorrigieren, jedoch nähren sich die Werte erstaunlich gut der $0.00 \sigma$-Umgebung an.
Dies liegt jedoch auch mitunter daran, dass die Fehler immer größer werden und diese immer schwerwiegender werden, da diese in der Qadratsumme stehen und deren Beitrag somit immer signifikanter wird, je größer die Ungenauigkeit wird. Dennoch bleiben die Werte verblüffend nah beieinander.


\paragraph{Reynodlszahl}
Zuletzt sollen noch die Reynodlszahlen betrachtet werden. Diese untereinander zu vergleichen würde keinen Sinn ergeben, daher wollen wir schauen, wie stark die Werte von den im Skript genannten kritischen Werten abweichen. \cite{skriptPAP21}

Somit kommt man für Stokes auf eine Abweichung von

\eq{abw_stokes}{
    \sigma_\text{Re, s} = 1.33\sigma \, .
}

Wie zu erwarten liegen wir ca. um den kritischen Punkt, denn dieser Punkt sollte gerade abgeschätzt werden. Da aber nun geschätzt wurde, hätte die Signifikanz bei einer guten Schätzung unter $1\sigma$ fallen müssen, daher lässt sich sagen, dass die Abschätzung etwas zu liberal war. Ab einem Fehler von $0.4$ wäre hier die Signifikanz unter $1$.

Für die Werte von Hagen-Poiseuille kommt man auf eine Signifikanz von über $1 \cdot 10^7 \sigma$. Dmait ist mehr als eindeutig, dass die Strömung während der Messung als Laminar angenommen werden konnte. 

Daraus lässt sich auch deuten, wieso die Methode nach Hagen-Poiseuille vermutlich präziser ist, als die von Stokes; denn die Annährung an die theoretisch laminare Strömung wird hier weitaus besser bewerkstelligt. 

\section{Kritik}
Zum Abschluss noch ein paar Kritiken. Als erstes wäre es besser gewesen, wenn der Volumenfluss über die Masse und Dichte des PEG bestimmt werden würde. Dann könnte man anstatt der visuellen Einschätzung eine Waage zur Hilfe nehmen und den Fehler minimieren. 
Außerdem wäre es vorteilhaft zu wissen, welches PEG es genau ist. Denn die PEG-Kettenlänge in der Industrie geht von gerade mal 100 Kettengliedern hin bis zu tausenden. Daher ist ein Vergleich mit online Werten witzlos. Auch ergibt die Grafik \afz{Dichte von Polyethylenglykol} erst richtig Sinn, wenn das genau PEG bekannt wäre.

\wrap[0.3]{L}{memes/ergebnisse}{Ergebnisse loben}
